\section{Feature Engineering Pipeline}
\label{sec:4_4}

The feature engineering pipeline transforms raw merged data into model-ready feature matrices, implemented in \texttt{src/features/schema.py} and \texttt{src/features/store.py}.

\subsection{Feature Schema Definition}
\label{sec:4_4_1}

The feature schema defines 65 base features across 5 categories:

\begin{itemize}
    \item \textbf{SEON Boolean} (21 features): \texttt{tor}, \texttt{vpn}, \texttt{public\_proxy}, \texttt{disposable\_email}, \texttt{phone\_is\_valid}, etc.
    \item \textbf{SEON Categorical} (12 features): \texttt{ip\_country}, \texttt{ip\_isp\_name}, \texttt{session/browser}, \texttt{session/device\_type}, etc.
    \item \textbf{Listing Numeric} (10 features): \texttt{LISTING\_PRICES\_RENT\_NET}, \texttt{LISTING\_PRICES\_BUY\_PRICE}, \texttt{LISTING\_CHARACTERISTICS\_ROOMS}, etc.
    \item \textbf{Listing Categorical} (7 features): \texttt{LISTING\_OFFERTYPE}, \texttt{LISTING\_CATEGORIES}, \texttt{TARGETPLATFORM}, etc.
    \item \textbf{Temporal} (4 features): \texttt{event\_hour}, \texttt{event\_day\_of\_week} (with cyclical encoding)
\end{itemize}

\subsection{Feature Preprocessing}
\label{sec:4_4_2}

Boolean features are cast to Int8 for XGBoost null handling. Categorical features use sklearn's LabelEncoder with unknown category handling, storing encoders for test-time application. Temporal features apply cyclical encoding to preserve periodicity: hour values (0-23) and weekday values (0-6) are transformed into sine/cosine pairs using the formula $\sin(2\pi x / period)$ and $\cos(2\pi x / period)$, ensuring that values at the boundary (e.g., 23:00 and 00:00) are numerically close in the feature space. Complete preprocessing implementation is provided in Appendix~A.

\textbf{Excluded Features:} STATUS column removed after leakage discovery; SEON ML scores excluded for model independence.

\subsection{Feature Schema Structure}
\label{sec:4_4_3}

The feature schema is implemented as a Python dataclass in \texttt{src/features/schema.py}, providing programmatic access to feature metadata. The \texttt{FeatureSchema} class maintains five lists corresponding to each feature category (seon\_boolean, seon\_categorical, listing\_numeric, listing\_categorical, temporal), with an \texttt{all\_features} property that concatenates these lists to produce the complete 65-feature set. This schema serves multiple purposes: (1) enforcing consistent feature selection across training and inference, (2) enabling automatic feature type detection for preprocessing, (3) facilitating feature importance analysis by category. Complete implementation is provided in Appendix~A.

\subsection{Feature Engineering Flow}

The complete feature engineering pipeline follows this sequence:

\begin{enumerate}
    \item \textbf{Data Loading:} \texttt{FeatureStore.load\_data()} reads the merged Parquet file with lazy evaluation
    \item \textbf{Datetime Parsing:} Convert timestamp columns to datetime64[ns] with UTC timezone
    \item \textbf{Temporal Feature Creation:}
    \begin{itemize}
        \item \texttt{event\_hour}: Extract hour (0-23)
        \item \texttt{event\_weekday}: Extract weekday (0-6, Monday=0)
        \item \texttt{event\_is\_weekend}: Binary flag for Saturday/Sunday
        \item \texttt{event\_is\_business\_hours}: Binary flag for 9am-5pm
    \end{itemize}
    \item \textbf{Boolean Casting:} Convert boolean columns to Int8 for XGBoost null handling
    \item \textbf{Categorical Encoding:} Apply LabelEncoder with unknown category support
    \item \textbf{Cyclical Encoding:} Transform hour and weekday to sin/cos pairs for periodicity
\end{enumerate}

\textbf{Model-Specific Feature Sets:}

\begin{table}[h]
\centering
\caption{Feature counts by model variant}
\begin{tabular}{lll}
\hline
\textbf{Model Variant} & \textbf{Base Features} & \textbf{GNN Embeddings} \\
\hline
Logistic Regression & 65 & 0 \\
Random Forest & 65 & 0 \\
Vanilla XGBoost & 65 & 0 \\
GNN+XGBoost & 65 & 16 \\
\hline
\textbf{Total} & & \textbf{81 (hybrid)} \\
\hline
\end{tabular}
\end{table}

For the hybrid GNN+XGBoost model, the 16-dimensional GraphSAGE embeddings are concatenated with the 65 base features after GNN training, resulting in an 81-dimensional feature vector for each listing node.
